\begin{lemma}[Tate]
\label{lemma-tate-japanese}
\begin{reference}
\cite[Theorem 23.1.3]{EGA}
\end{reference}
Let $R$ be a ring.
Let $x \in R$.
Assume
\begin{enumerate}
\item $R$ is a normal Noetherian domain,
\item $R/xR$ is a domain and N-2,
\item $R \cong \lim_n R/x^nR$ is complete with respect to $x$.
\end{enumerate}
Then $R$ is N-2.
\end{lemma}

\begin{proof}
We may assume $x \not = 0$ since otherwise the lemma is trivial.
Let $K$ be the fraction field of $R$. If the characteristic of $K$
is zero the lemma follows from (1), see
Lemma \ref{lemma-domain-char-zero-N-1-2}. Hence we may assume
that the characteristic of $K$ is $p > 0$, and we may apply
Lemma \ref{lemma-domain-char-p-N-1-2}. Thus given $L/K$
a finite purely inseparable field extension we have to show
that the integral closure $S$ of $R$ in $L$ is finite over $R$.

\medskip\noindent
Let $q$ be a power of $p$ such that $L^q \subset K$.
By enlarging $L$ if necessary we may assume there exists
an element $y \in L$ such that $y^q = x$. Since $R \to S$
induces a homeomorphism of spectra (see Lemma \ref{lemma-p-ring-map})
there is a unique prime ideal $\mathfrak q \subset S$ lying
over the prime ideal $\mathfrak p = xR$. It is clear that
$$
\mathfrak q = \{f \in S \mid f^q \in \mathfrak p\} = yS
$$
since $y^q = x$. Observe that $R_{\mathfrak p}$ is a discrete
valuation ring by Lemma \ref{lemma-characterize-dvr}. Then
$S_{\mathfrak q}$ is Noetherian by Krull-Akizuki
(Lemma \ref{lemma-krull-akizuki}). Whereupon we conclude
$S_{\mathfrak q}$ is a discrete valuation ring by
Lemma \ref{lemma-characterize-dvr} once again.
By Lemma \ref{lemma-finite-extension-residue-fields-dimension-1} we
see that $\kappa(\mathfrak q)/\kappa(\mathfrak p)$ is
a finite field extension. Hence the integral closure
$S' \subset \kappa(\mathfrak q)$ of $R/xR$ is finite over
$R/xR$ by assumption (2). Since $S/yS \subset S'$ this implies
that $S/yS$ is finite over $R$. Note that $S/y^nS$ has a finite
filtration whose subquotients are the modules
$y^iS/y^{i + 1}S \cong S/yS$. Hence we see that each $S/y^nS$
is finite over $R$. In particular $S/xS$ is finite over $R$.
Also, it is clear that $\bigcap x^nS = (0)$ since an element
in the intersection has $q$th power contained in $\bigcap x^nR = (0)$
(Lemma \ref{lemma-intersect-powers-ideal-module-zero}).
Thus we may apply Lemma \ref{lemma-finite-over-complete-ring} to conclude
that $S$ is finite over $R$, and we win.
\end{proof}